\chapter{Pierre René Deligne}

\addcontentsline{toc}{chapter}{Problem History}
\section*{Problem History}

\begin{prereqs}
\noindent \textbf{Background}
\begin{itemize}
\item Algebraic geometry of curves over finite fields
\item Cohomology theory for algebraic varieties
\item zeta functions and L-functions
\end{itemize}
\end{prereqs}

\begin{objectives}
\noindent \textbf{Learning Objectives}
\begin{itemize}
\item Understand the Weil conjectures and their origin in the study of algebraic varieties over finite fields
\item Learn how étale cohomology provides a cohomological framework for the Lefschetz fixed-point formula
\item Appreciate Deligne's proof as the culmination of decades of work by Weil, Grothendieck, and others
\end{itemize}
\end{objectives}

The Weil conjectures, formulated by André Weil in 1949, describe the
zeta function of a smooth projective variety over a finite field,
predicting rationality, a functional equation, and---most
remarkably---that zeros and poles have geometric origin.

Given a smooth projective variety $X$ over $\mathbb{F}_q$, the zeta
function is
\[
Z(X, t) = \exp\!\Bigg(\sum_{n=1}^{\infty} \frac{|X(\mathbb{F}_{q^n})|}{n}\, t^n\Bigg).
\]
Rationality was proved by Dwork (1960) using $p$-adic methods, but the
deeper geometric meaning required a cohomology theory for varieties over
arbitrary fields.

Grothendieck's insight was that \'etale cohomology---replacing the
analytic topology with the \'etale topology of scheme morphisms---should
yield a Lefschetz trace formula for Frobenius endomorphisms. The
resulting cohomology groups $H^i_{\text{ét}}(X, \mathbb{Q}_\ell)$ satisfy
axioms paralleling singular cohomology, and the trace formula gives
\[
|X(\mathbb{F}_{q^n})| = \sum_{i=0}^{2d} (-1)^i \operatorname{tr}\!\big(\text{Frob}_q^n \mid H^i_{\text{ét}}(X, \mathbb{Q}_\ell)\big).
\]
Rationality and the functional equation follow. The Riemann hypothesis:
eigenvalues of $\text{Frob}_q$ on $H^i$ have absolute value $q^{i/2}$.
Deligne announced a proof in 1973, published in 1974 \cite{deligne1974}.

\addcontentsline{toc}{chapter}{How Was It Solved}
\section*{How Was It Solved}

The key challenge is controlling Frobenius eigenvalues on intermediate
cohomology $H^i$ for $2 \leq i \leq 2d - 2$. Deligne's approach uses
algebraic cycles and bilinear pairings.

The Hard Lefschetz theorem (proved by Deligne using Hironaka's resolution
of singularities) states that for an ample divisor $H$, the cup product
\[
L^k: H^{d-k}_{\text{ét}}(X, \mathbb{Q}_\ell) \longrightarrow H^{d+k}_{\text{ét}}(X, \mathbb{Q}_\ell)
\]
is an isomorphism, where $L$ is the class of $H$.

The critical innovation is the Hodge-Riemann bilinear relations on
primitive cohomology
$P^k = \ker(L^{d-k+1}: H^k \to H^{2d-k+2})$. Deligne proves definite
signature conditions---an analogue of Hodge-Riemann relations from
K\"{a}hler geometry---via a limiting argument showing the signature is a
topological invariant.

To bridge algebraic de Rham and \'etale cohomology, Deligne invokes
Grothendieck's comparison theorems and the degeneration of the
Hodge-de Rham spectral sequence. The Hodge-Riemann relations force the
eigenvalues to have correct absolute values.

Deligne's second proof (1980) streamlined the argument by working
directly with the weight filtration on \'etale cohomology, removing the
need for Hard Lefschetz. This version has had broader applications.

\addcontentsline{toc}{chapter}{Impact}
\section*{Impact}

Deligne's proof confirmed Grothendieck's vision that \'etale cohomology
is the correct framework for arithmetic geometry, validating scheme
theory and motivic cohomology.

The theory of weights became a cornerstone of modern algebraic geometry,
organizing geometric information precisely. It influences moduli spaces,
the Langlands program, and Shimura varieties. The Weil conjectures provide
the foundation for L-functions of algebraic varieties, with the Riemann
hypothesis giving powerful point-counting estimates.

The techniques found applications through $\ell$-adic representations:
\'etale cohomology carries absolute Galois group action encoding deep
arithmetic information. Deligne's monodromy theorem---showing the
monodromy group of a smooth projective family is semisimple---has found
many applications.

\addcontentsline{toc}{chapter}{Modern Views}
\section*{Modern Views}

The Langlands program uses weights and L-functions as building blocks.
The Sato-Tate conjecture (proved by Taylor et al.) concerns Frobenius
eigenvalue distributions for elliptic curves.

The theory of perverse sheaves, growing from Deligne's work, is central to
geometric representation theory. The decomposition theorem of
Beilinson--Bernstein--Deligne--Gabber---that proper map cohomology splits
as intersection cohomology sheaves---is one of the most powerful tools in
the field.

The Weil conjectures motivated $p$-adic Hodge theory, comparing
cohomology theories in mixed characteristic via crystalline cohomology
and the de Rham-Witt complex. In recent decades, they have found
applications in combinatorics and theoretical physics.

Deligne received the Fields Medal in 1978 for this work, widely regarded
as one of the greatest mathematical achievements of the twentieth century.

\section*{Summary}
The Weil conjectures describe the zeta function of a smooth projective
variety over a finite field. Deligne proved the Riemann hypothesis using
\'etale cohomology and geometric analysis of algebraic cycles, confirming
Grothendieck's vision and establishing the theory of weights as a
fundamental tool in algebraic geometry.

\addcontentsline{toc}{chapter}{Exercises}
